% FAC-EN-U0054
% Source: fac_body.tex lines 1927--1952.
% Authority: PDF physical pp. 53--54; printed pp. 248--249 (no. 56 complete).
% English source-first translation; mathematical notation and definitions preserved.

\phantomsection\label{fac:III:2:no:56}
\textbf{56. Graded modules.}

\phantomsection\label{fac:III:2:no:56:def:graded-algebra}
Let $S=K[t_0,\ldots,t_r]$ be the polynomial algebra in
$t_0,\ldots,t_r$. For every integer $n\geq0$, let $S_n$ be the vector
subspace of $S$ formed by the homogeneous polynomials of degree $n$; for
$n<0$, put $S_n=0$. The algebra $S$ is the direct sum of the $S_n$,
$n\in\mathbf{Z}$, and
\[
S_pS_q\subset S_{p+q};
\]
in other words, $S$ is a \textit{graded algebra}.

\phantomsection\label{fac:III:2:no:56:def:graded-module}
Recall that an $S$-module $M$ is called \textit{graded} when a decomposition
of $M$ as a direct sum
\[
M=\sum_{n\in\mathbf{Z}}M_n
\]
has been specified, where the $M_n$ are subgroups of $M$ such that
$S_pM_q\subset M_{p+q}$ for every pair of integers $(p,q)$. An element of
$M_n$ is called \textit{homogeneous} of degree $n$. A submodule $N$ of $M$
is called \textit{homogeneous} if it is the direct sum of the $N\cap M_n$,
in which case it is a graded $S$-module.

\phantomsection\label{fac:III:2:no:56:def:homogeneous-homomorphism}
If $M$ and $M'$ are two graded $S$-modules, an $S$-homomorphism
\[
\varphi:M\longrightarrow M'
\]
is called \textit{homogeneous of degree $s$} if
$\varphi(M_n)\subset M'_{n+s}$ for every $n\in\mathbf{Z}$. A homogeneous
$S$-homomorphism of degree zero will simply be called a
\textit{homomorphism}.

\phantomsection\label{fac:III:2:no:56:def:shift}
If $M$ is a graded $S$-module and $n$ is an integer, we denote by $M(n)$
the graded $S$-module
\[
M(n)=\sum_{p\in\mathbf{Z}}M(n)_p,
\qquad M(n)_p=M_{n+p}.
\]
Thus $M(n)=M$ as an $S$-module, but an element homogeneous of degree $p$ in
$M(n)$ is homogeneous of degree $n+p$ in $M$; in other words, $M(n)$ is
obtained from $M$ by lowering the degrees by $n$ units.

\phantomsection\label{fac:III:2:no:56:def:class-C}
We denote by $\mathcal{C}$ the class of graded $S$-modules $M$ such that
$M_n=0$ for every sufficiently large $n$. If
\[
A\longrightarrow B\longrightarrow C
\]
is an exact sequence of homomorphisms of graded $S$-modules, the relations
$A\in\mathcal{C}$ and $C\in\mathcal{C}$ clearly imply
$B\in\mathcal{C}$; in other words, $\mathcal{C}$ is indeed a class in the
sense of [14], Chapter~I.

\phantomsection\label{fac:III:2:no:56:def:C-morphism}
More generally, we shall use the terminology introduced in the cited
article. In particular, a homomorphism $\varphi:A\to B$ is called
$\mathcal{C}$-\textit{injective} (respectively
$\mathcal{C}$-\textit{surjective}) if
$\operatorname{Ker}(\varphi)\in\mathcal{C}$ (respectively
$\operatorname{Coker}(\varphi)\in\mathcal{C}$), and
$\mathcal{C}$-\textit{bijective} if it is both $\mathcal{C}$-injective and
$\mathcal{C}$-surjective.

\phantomsection\label{fac:III:2:no:56:def:TF}
A graded $S$-module $M$ is called \textit{finitely generated} if it is
generated by finitely many elements. We say that $M$ \textit{satisfies condition (TF)}
if there is an integer $p$ such that the submodule
\[
\sum_{n\geq p}M_n
\]
of $M$ is finitely generated. Equivalently, $M$ is
$\mathcal{C}$-isomorphic to a finitely generated module. The modules
satisfying (TF) form a class containing $\mathcal{C}$.

\phantomsection\label{fac:III:2:no:56:def:free-graded-module}
A graded $S$-module $L$ is called \textit{free} (respectively
\textit{finitely generated free}) if it has a basis (respectively a finite
basis) formed by homogeneous elements; in other words, if it is isomorphic
to a direct sum (respectively a finite direct sum) of modules $S(n_i)$.
